\chapter{2005:Yves Chauvin and Robert H. Grubbs and Richard R. Schrock}

\label{chap:chemistry2005}

The Nobel Prize in Chemistry for the year 2005 was awarded jointly to Yves Chauvin and Robert H. Grubbs and Richard R. Schrock.

\textbf{Motivation:}

for the development of the metathesis method in organic synthesis

\section{Yves Chauvin}

\subsection*{Early Life and Education}

Yves Chauvin was born in 1930 in Menen, Belgium.

\subsection*{Career and Major Research}

Chauvin studied at the École supérieure de chimie, physique, électronique de Lyon and spent his career at the Institut Français du Pétrole (IFP) in Rueil-Malmaison, France, where he became honorary research director. In 1971 he proposed the mechanism of olefin metathesis, in which a metal carbene and a metallacyclobutane intermediate carry the breaking and re-forming of carbon--carbon double bonds.

\subsection*{Nobel Prize-winning Work}

The work for which Yves Chauvin was awarded the Nobel Prize in Chemistry in 2005 was described by the Nobel Committee as: "for the development of the metathesis method in organic synthesis".

Chauvin's mechanism explained how the double bonds of alkenes are broken and reorganized in the presence of certain metal complexes, giving the reaction a rational basis that other chemists could build upon.

\subsection*{Significance and Impact}

By explaining the metathesis mechanism, Chauvin's work enabled the later development of practical catalysts that made metathesis a standard reaction in organic synthesis. The reaction is now used industrially to make polymers, fine chemicals, and pharmaceuticals.

\subsection*{Later Career and Honors}

Chauvin retired from the Institut Français du Pétrole and was elected to the French Academy of Sciences. He died on 2015-01-27 in Tours, France.

\subsection*{Legacy}

Chauvin's mechanistic insight underlies every modern metathesis catalyst, including the ruthenium and molybdenum systems recognized in the same prize.

\subsection*{Modern Developments}

Chauvin's elucidation of the mechanism of olefin metathesis, based on metallacyclobutane intermediates, provided the theoretical foundation for the metathesis reaction. His mechanism guided the development of the Grubbs and Schrock catalysts that now enable the efficient synthesis of pharmaceuticals, polymers, and fine chemicals.

\subsection*{Selected Publications}

"Catalyse de Transformation des Oléfines par les Complexes du Tungstène" (Die Makromolekulare Chemie, 1971)

\clearpage

\section{Robert H. Grubbs}

\subsection*{Early Life and Education}

Robert H. Grubbs was born in 1942 in Possum Trot, Kentucky, USA.

\subsection*{Career and Major Research}

Grubbs earned a bachelor's degree from the University of Florida and a doctorate from Columbia University in 1968, followed by postdoctoral work at Stanford University. He became a professor at the California Institute of Technology, where in the early 1990s he discovered ruthenium-based metathesis catalysts that are stable in air and tolerant of many functional groups.

\subsection*{Nobel Prize-winning Work}

The work for which Robert H. Grubbs was awarded the Nobel Prize in Chemistry in 2005 was described by the Nobel Committee as: "for the development of the metathesis method in organic synthesis".

The Grubbs catalysts brought olefin metathesis out of the specialist laboratory, making the reaction easy to run under ordinary conditions.

\subsection*{Significance and Impact}

Because the ruthenium catalysts work in air and tolerate a wide range of chemical groups, they allowed metathesis to be applied broadly in organic synthesis. The reaction is now used to produce pharmaceuticals, polymers, and agrochemicals.

\subsection*{Later Career and Honors}

Grubbs remained professor of chemistry at the California Institute of Technology for the rest of his career. He died on 2021-12-19 in Duarte, California, USA.

\subsection*{Legacy}

The Grubbs catalyst is one of the most widely used catalysts in synthetic chemistry and is standard equipment in research and industry.

\subsection*{Modern Developments}

Grubbs's metathesis catalysts, including the ruthenium-based Grubbs catalyst, made olefin metathesis a practical, widely used reaction in organic synthesis. Metathesis is now used in the industrial production of pharmaceuticals, agrochemicals, and specialty polymers, and his catalysts transformed the field of green chemistry.

\subsection*{Selected Publications}

"The Development of L2X2Ru=CHR Olefin Metathesis Catalysts" (Accounts of Chemical Research, 2001)

\clearpage

\section{Richard R. Schrock}

\subsection*{Early Life and Education}

Richard R. Schrock was born in 1945 in Berne, Indiana, USA.

\subsection*{Career and Major Research}

Schrock earned bachelor's and doctoral degrees from Harvard University, in 1967 and 1971, and completed postdoctoral work at the University of Cambridge. After a period at DuPont Central Research he joined the faculty of the Massachusetts Institute of Technology in 1975. He developed highly active molybdenum and tungsten alkylidene metathesis catalysts.

\subsection*{Nobel Prize-winning Work}

The work for which Richard R. Schrock was awarded the Nobel Prize in Chemistry in 2005 was described by the Nobel Committee as: "for the development of the metathesis method in organic synthesis".

Schrock's molybdenum alkylidene catalysts were the first to show exceptional activity in olefin metathesis, revealing the power of the reaction.

\subsection*{Significance and Impact}

Schrock's catalysts made it possible to conduct metathesis efficiently at low temperatures and to control the products precisely. They remain important for reactions such as ring-opening metathesis polymerization.

\subsection*{Later Career and Honors}

Schrock remains a professor of chemistry at the Massachusetts Institute of Technology, where he continues research on organometallic chemistry and catalysis.

\subsection*{Legacy}

Together with the catalysts of Grubbs, Schrock's alkylidene systems made olefin metathesis one of the central reactions of modern organic synthesis.

\subsection*{Modern Developments}

Schrock's molybdenum and tungsten metathesis catalysts, the first to achieve controlled olefin metathesis, transformed the reaction from a curiosity into a powerful synthetic tool. Metathesis is now used to make pharmaceuticals, advanced polymers, and fuels, and Schrock's work pioneered the field of high-oxidation-state organometallic catalysis.

\subsection*{Selected Publications}

"Preparation and Characterization of Molybdenum and Tungsten Alkylidene Complexes" (Journal of the American Chemical Society, 1990)

\clearpage

\section*{Chapter Summary}

The 2005 prize recognized the three chemists who turned olefin metathesis into a practical tool of organic synthesis: Chauvin's mechanism, Schrock's molybdenum alkylidene catalysts, and Grubbs's air-stable ruthenium catalysts. Metathesis is now used to make polymers, agrochemicals, and pharmaceuticals.
